% !TEX root = main.tex
\documentclass{article}
\usepackage{graphicx} % Required for inserting images
\usepackage[margin=2.5cm]{geometry}
\usepackage[expansion=false]{microtype}
\usepackage{titlesec}
\sloppy
\titleformat{\section}{\normalfont\Large\bfseries\centering}{}{0em}{}

\title{PE\&RC PhD Project Proposal}
\author{Tom Brouwers}
\date{March 2026}

\begin{document}

\maketitle

\section{General Information}
\subsection{About the Project and Evaluation Criteria:}
The PE\&RC PhD proposal is written by the PhD candidate in consultation with and feedback from the supervisory team. After submission, the proposal is sent out for review. The candidate and the supervisory team are encouraged to suggest possible independent reviewers that can provide constructive feedback to the candidate. 

The reviewers are asked to evaluate the proposal on:
\begin{itemize}
\item Scientific quality
Clarity of writing, depth of literature review, definition of research question(s) and objective(s), and originality of subject and/or methodology.
\item Feasibility
Feasibility given the 4-year time frame, (financial) resources, effectiveness of the proposed approach/strategy, logistical/practical aspects and evaluation of the expertise within the supervising team and/or collaborations for adequate support in the proposed project.
\item Scientific and/or societal impact
The importance of the potential research results in the short and long term in the field of research, potential use and relevance in other scientific disciplines, and potential use and relevance for society (economic, technical, social, cultural), including technological application. 
\end{itemize}
\subsection{General Project Information: }
\begin{table}[h]
    \centering
    \begin{tabular}{|p{8cm}|p{8cm}|}
    \hline
    Main PE\&RC affiliated Institute/University  & Wageningen University \\
       \hline
      Main PE\&RC research group & Soil, Geography and Landscape \\
      \hline
      Other PE\&RC research groups involved & \\
      \hline
      Where will the research be conducted (country) & The Netherlands \\
      \hline
      Full name of the PhD candidate & Thomas Fidélis Brouwers \\
      \hline
      Project duration & FROM 01/01/2026 TO 01/01/2030 \\
      \hline
      Time spent on PhD project & \textbullet\ Fulltime \\
      \hline
      Principal supervisor/Promotor & Prof. Veerle Vanacker \\
      \hline
      Additional supervisor(s) & Dr. Annegret Larsen; Prof. Benjamin Vernot; Dr. Peter van der Sleen \\
      \hline
    \end{tabular}
    \label{tab:general_info}
\end{table}
\subsection{Ethics:}
Ethical issues such as experiments with animals (vertebrates), work with genetically modified organisms, work with material under the Nagoya protocol, or personal data privacy considerations (AVG) may have legal implications. Also, ethical guidelines of the respective research institution may apply.
\begin{table}[h]
    \centering
    \begin{tabular}{|p{12cm}|p{4cm}|}
        \hline
        Do any of the above mentioned, or other ethical considerations apply to this project? & NO \\
         \hline
         \multicolumn{2}{|p{16cm}|}{if YES, please elaborate: } \\ 
         \hline
    \end{tabular}
    \label{ethics}
\end{table}
\subsection{PE\&RC Research Theme(s)}
In which of the PE\&RC research theme(s) does the project fit (multiple options possible). Please elaborate.
\begin{table}[h]
    \centering
    \begin{tabular}{|p{8cm}|p{8cm}|}
        \hline
        \begin{minipage}[t]{7.8cm}
            \textbullet\ Biodiversity and System Dynamics\\
            \textbullet\ Redesign of (agro-) ecological systems \& landscapes
        \end{minipage}
        &  
        The project aims to better understand the biological composition of ancient, natural rivers in the European context, and through this re-evaluate the role of biotic drivers and ecosystem engineers which are crucial for maintaining the form and function of said ecosystems. These insights will be used to inform river restoration projects and, hopefully, lead to higher success rates in creating and maintaining resilient and biodiverse river ecosystems. To achieve these goals the method of sedimentary ancient DNA analysis will be further developed and then applied in various fluvial context across Europe.\\
        \hline
    \end{tabular}
    \label{themes}
\end{table}
\section{Proposal Structure}
\subsection{Title }
\begin{center}
    \textbf{Sedimentary Ancient DNA as a Robust Tool for River Science}
\end{center}

\subsection{Abstract (max. 250 words): }
Rivers restoration is crucial, strives for restoring rivers to 'natural' systems, but baseline is missing. This study aims to develop the method of sedimentary ancient DNA (sedaDNA) analysis for river science, to reconstruct the biological composition of ancient, natural rivers in the European context. 
\subsection{Research Topic and Background}




\bibliography{references}
\bibliographystyle{abbrv}
\end{document}
